\documentclass[aps,pra,onecolumn,superscriptaddress,nofootinbib,longbibliography]{revtex4-2}

\usepackage[T1]{fontenc}
\usepackage{lmodern}
\usepackage{amsmath}
\usepackage{amssymb}
\usepackage{amsthm}
\usepackage{booktabs}
\usepackage[colorlinks=true,linkcolor=blue,citecolor=blue,urlcolor=blue]{hyperref}

\theoremstyle{plain}
\newtheorem{theorem}{Theorem}
\newtheorem{lemma}[theorem]{Lemma}
\newtheorem{proposition}[theorem]{Proposition}
\newtheorem{corollary}[theorem]{Corollary}
\theoremstyle{remark}
\newtheorem{remark}[theorem]{Remark}

\newcommand{\F}{\mathrm{F}}
\newcommand{\LC}{\mathrm{LC}}
\newcommand{\rw}{\mathrm{rw}}
\newcommand{\DHclass}{\mathsf{DH}}
\newcommand{\Prod}{\mathsf{Prod}}
\newcommand{\codeus}{\textunderscore\allowbreak}
\newcommand{\code}[1]{\texttt{\let\_\codeus #1}}
\newcommand{\nbr}{N}

\begin{document}

\title{Distance-hereditary graphs characterize minimum-fusion synthesis of graph states}

\author{Sean Howell}
\email{sean.g.howell@gmail.com}

\begin{abstract}
In fusion-based quantum computation, graph states are assembled from small entangled resource
states by destructive two-qubit measurements called fusions. We study the minimum number $\F(G)$ of
fusions needed to produce the graph state of a graph $G$ on $n$ vertices from three-qubit
Greenberger-Horne-Zeilinger states, with Bell-type fusions on arbitrary pairs of qubits and free
single-qubit Clifford operations. We prove that $\F(G)\ge n+3\,\rw(G)-6$ for every connected graph
on $n\ge 3$ vertices, where $\rw(G)$ is the rank-width. The proof turns a spanning tree of the
network formed by the resources and the fusions into a rank decomposition of $G$ and bounds the Schmidt
rank of the output across each of its cuts in terms of the cycle rank of the network. The bound is
attained at rank-width one, so that for $n\ge 3$ we have $\F(G)=n-3$ if and only if $G$ is
distance-hereditary. Both directions of this characterization can be assembled from published results
on Schmidt-rank width and split decompositions, although we have not found it stated in the literature.
The graph-theoretic core of its constructive direction is verified in Lean 4. An exhaustive
search over the 585 connected local-complementation orbits on three to nine vertices, with every
construction checked by two independently implemented stabilizer simulators, gives nineteen exact
values of $\F$ above the minimum. The sixteen orbits with $\F=n$ attain the rank-width bound, whereas
three orbits of rank-width two require $n+3$ fusions, so above rank-width one the fusion cost is not
determined by rank-width. For 25 orbits of rank-width three whose minimum is not known, we improve the
lower bound from $n$ to $n+3$.
\end{abstract}

\maketitle

\section{Introduction}\label{sec:intro}

Fusion-based quantum computation (FBQC) builds fault-tolerant computation from a supply of small,
identical entangled resource states that are joined by fusions, destructive joint measurements of two
qubits~\cite{Bartolucci2023fusion}. The three-qubit Greenberger-Horne-Zeilinger (GHZ) state is a basic
resource in several photonic proposals~\cite{GimenoSegovia2015three,Zaidi2015near}, and the number of
fusions is a main cost in such architectures, which raises the question of how few fusions suffice to
prepare a given entangled state. For graph states, the stabilizer states in one-to-one correspondence
with simple graphs~\cite{Hein2004multiparty,Hein2006entanglement}, this becomes a combinatorial
optimization over graphs. It has been approached with graph-theoretical heuristics for GHZ resources~\cite{Lee2023graph},
with dynamic programming over local-complementation orbits for caterpillar states produced by a single
quantum emitter~\cite{Lobl2025generating}, and through the design of fusion networks built from linear
cluster states~\cite{Staudacher2026adaptive,Cashman2025finding} or from ring, star and linear motifs joined
by type-I fusion~\cite{Ji2026towards}.

We take the resource to be the three-qubit GHZ state, whose graph is the path $P_3$. A fusion is the
destructive measurement of a commuting pair of two-qubit Pauli operators that distinguishes the four Bell
states, and it may act on any two qubits. Single-qubit Clifford operations are free. We write $\F(G)$ for
the least number of fusions in a schedule whose output equals the graph state $|G\rangle$ up to
single-qubit Cliffords and a relabeling of the qubits. Counting qubits and connected components shows
that $\F(G)\ge n-3$ for every connected graph on $n\ge 3$ vertices (Proposition~\ref{prop:count}). We
prove a stronger lower bound in terms of the rank-width $\rw(G)$ and characterize the graphs that attain
$n-3$.

\begin{theorem}[Rank-width bound]\label{thm:rw}
For every connected graph $G$ on $n\ge 3$ vertices, $\F(G)\ge n+3\,\rw(G)-6$.
\end{theorem}

\begin{theorem}[Characterization]\label{thm:main}
For a connected graph $G$ on $n\ge 3$ vertices, $\F(G)=n-3$ if and only if $G$ is
distance-hereditary.
\end{theorem}

Distance-hereditary graphs~\cite{Howorka1977characterization,Bandelt1986distance} are the graphs in
which every connected induced subgraph preserves distances. They coincide with the graphs of rank-width
at most one~\cite{Oum2005rank}, with the graphs that have no vertex-minor isomorphic to the
five-cycle~\cite{Bouchet1987isotropic,Bouchet1988transforming,Oum2017rank}, and with the graphs obtained
from a single vertex by repeatedly adding pendant vertices and twins~\cite{Bandelt1986distance}. Since
rank-width is invariant under local complementation~\cite{Oum2005rank} and equals the Schmidt-rank width
of the graph state~\cite{VandenNest2006universal,VandenNest2007classical}, both theorems relate the
fusion cost to an entanglement measure of $|G\rangle$ that is invariant under local Clifford operations.

Kumabe, Mori and Yoshimura~\cite{Kumabe2026complexity} proved the analogous results for Clifford
circuits, where preparing a connected graph state requires at least $n+\rw(G)-2$ two-qubit gates, and
exactly $n-1$ when the rank-width is one. For caterpillar states produced by a single quantum emitter,
L\"obl \emph{et al.}~\cite{Lobl2025generating} computed the exact minimum number of fusions for all graph
states on up to eight qubits and showed that it is at least the linear rank-width minus one. The graphs
that need no fusion in their model are those locally equivalent to caterpillars, which are the connected
graphs of linear rank-width one~\cite{BuiXuan2013note}, and both the minimum fusion counts and the class
of graphs at the minimum differ from ours. Theorem~\ref{thm:rw} is the counterpart of these two bounds for
fusions of three-qubit GHZ states.
Both directions of Theorem~\ref{thm:main} can also be assembled from published results. The implication
from $\F(G)=n-3$ to distance-heredity follows from the equality of Schmidt-rank width and
rank-width~\cite{VandenNest2007classical} once one observes that the fusion network of a schedule with
$n-3$ fusions is a tree. The converse follows from the pendant-and-twin construction of Bandelt and
Mulder~\cite{Bandelt1986distance} together with the graph rule for fusion~\cite{Zaidi2015near,Lee2023graph},
or from the construction of distance-hereditary graph states by fusions along their split decomposition
by Connolly \emph{et al.}~\cite{Connolly2026efficient}, combined with the preparation of star graph states
from GHZ states by Lee and Jeong~\cite{Lee2023graph}. We have not found Theorem~\ref{thm:main}, the bound
of Theorem~\ref{thm:rw}, or exact minimum fusion counts for this model in the literature. We make the
characterization explicit for this resource model and prove the more general bound.

Other preparation costs have been studied in related settings. Davies and Jena~\cite{Davies2025preparing}
bound the number of two-qubit gates needed to prepare graph states from classes closed under
vertex-minors, unconditionally for circle graphs and for bounded rank-width, and by $O(n\log n)$ for every
proper vertex-minor-closed class assuming a structural conjecture of Geelen. Cabello
\emph{et al.}~\cite{Cabello2011optimal} and Sharma \emph{et al.}~\cite{Sharma2026minimising} studied the
representatives of local-complementation orbits with the fewest gates or edges, and Sharma \emph{et al.}
used edge minimization to reduce the GHZ resources and fusions needed for repeater graph states. Adcock
\emph{et al.}~\cite{Adcock2019hard} bounded which graph states can be produced from optical resources
with postselected fusions. Whether one graph state can be transformed into another by single-qubit
Clifford operations, single-qubit Pauli measurements and classical communication is decidable in time
cubic in the size of the source when its rank-width and the size of the target are
bounded~\cite{Dahlberg2018transforming}. In general this problem is NP-complete, even when the target is a
tensor product of Bell pairs on specified pairs of qubits, and more efficient algorithms exist for
particular targets~\cite{Dahlberg2020how,Dahlberg2020transforming}. Counting the labeled graph states in
an orbit under local complementation alone, without identifying isomorphic graphs, is
\#P-complete~\cite{Dahlberg2020counting}. For distance-hereditary graphs, Connolly, Nishio and
Nemoto~\cite{Connolly2026local} derived the sizes of the local-equivalence classes of several families.
The local-complementation orbits of graph states were classified by Hein
\emph{et al.}~\cite{Hein2004multiparty} for $n\le 7$ and by Danielsen and
Parker~\cite{Danielsen2006classification} for $n\le 12$, and they were mapped by Adcock
\emph{et al.}~\cite{Adcock2020mapping}.

Theorems~\ref{thm:rw} and~\ref{thm:main} are proved in Secs.~\ref{sec:model} and~\ref{sec:main}. The proofs
use the graph rule for a fusion of two disjoint graph states (Lemma~\ref{lem:merge}) and standard facts
about rank-width. The graph-theoretic core of the constructive direction of Theorem~\ref{thm:main} has
been verified with the Lean~4 proof assistant~\cite{deMoura2021lean} and its mathematical
library~\cite{Mathlib2020lean}. The Lean development shows that a graph on $k+3$ vertices is reachable
from $P_3$ by $k$ pendant and false-twin additions interleaved with local complementations if and only if
it belongs to the class generated from the connected three-vertex graphs by pendant and twin additions and
local complementation (Sec.~\ref{sec:lean}). By results of Bandelt and Mulder~\cite{Bandelt1986distance}
and of Oum~\cite{Oum2005rank}, which are not part of the formalization, this class consists of the
connected distance-hereditary graphs on at least three vertices. In Sec.~\ref{sec:table} we report an exhaustive search over the $585$
connected local-complementation orbits with $3\le n\le 9$, which agrees with Theorem~\ref{thm:main} on
every orbit and determines nineteen values of $\F$ above the minimum. The search was run with an
automated system that accepts a claim only after machine verification. The coding assistant used to build
the system first noticed the characterization in the table of all $585$ orbits, and a language model
working within the system then re-derived it from the labeled orbits with $3\le n\le 7$ alone and wrote
the Lean proofs. A specialist can readily derive the characterization from the results cited above, and
we describe how it was found in Sec.~\ref{sec:method}.

\section{Graph states and fusions}\label{sec:model}

For a simple graph $G=(V,E)$ the graph state
\[
  |G\rangle=\prod_{\{u,v\}\in E}\mathrm{CZ}_{uv}\,|{+}\rangle^{\otimes V}
\]
is the unique common $+1$ eigenvector of the operators $K_v=X_v\prod_{u\in\nbr(v)}Z_u$, $v\in V$, where
$\nbr(v)$ is the neighborhood of $v$. Local complementation at $v$ replaces the subgraph induced on
$\nbr(v)$ by its complement, and we write $\tau_v(G)$ for the result. On graph states it is implemented,
up to a global phase, by a product of single-qubit Clifford unitaries~\cite{VandenNest2004graphical},
\begin{equation}\label{eq:lc}
  |\tau_v(G)\rangle = e^{-i\pi X_v/4}\prod_{u\in\nbr(v)}e^{i\pi Z_u/4}\,|G\rangle,
\end{equation}
and two graph states on the same qubits are related by a product of single-qubit Clifford unitaries if
and only if their graphs are related by a sequence of local complementations~\cite{VandenNest2004graphical}.
The qubits that survive a schedule carry no fixed labels, so we also identify isomorphic graphs and write
$G\sim_{\LC}H$ for the equivalence generated by local complementation and isomorphism. For
$A\subseteq V$, the cut-rank $\rho_G(A)$ is the rank over $\mathrm{GF}(2)$ of the $A\times(V\setminus A)$
submatrix of the adjacency matrix, and the Schmidt rank of $|G\rangle$ across the bipartition
$(A,V\setminus A)$ is $2^{\rho_G(A)}$~\cite{Hein2004multiparty}.

A schedule starts from $g$ copies of the three-qubit GHZ state, whose graph is $P_3$ and which is
$\LC$-equivalent to the triangle $K_3$. A fusion on two distinct qubits $a$ and $b$ is the destructive
measurement of the commuting observables $X_aZ_b$ and $Z_aX_b$, after which both qubits are discarded.
This is a Bell measurement preceded by a Hadamard gate on $b$, the convention of
Refs.~\cite{Lee2023graph,Lobl2025transforming}, and other Bell-type conventions differ from it by
single-qubit Cliffords on $a$ and $b$, which are free in our model. A schedule interleaves fusions, on any
pairs of qubits, with arbitrary single-qubit Clifford gates, and it produces $G$ if for every measurement
outcome the state of the surviving qubits is $\LC$-equivalent to $|G\rangle$. The outcomes of a fusion on a
stabilizer state give post-measurement states that differ by Pauli
operators~\cite{Lee2023graph,Lobl2025transforming}, so this condition does not depend on the outcomes. We
therefore treat fusions as deterministic and do not model the probabilistic success of linear-optical
fusion~\cite{Browne2005resource}. The cost of a schedule is its number
of fusions, and $\F(G)$ is the least cost of a schedule that produces $G$.

\begin{lemma}[Fusion of disjoint graph states]\label{lem:merge}
Let $G_1$ and $G_2$ be graphs on disjoint vertex sets, $a$ a vertex of $G_1$ and $b$ a vertex of $G_2$.
A fusion on $(a,b)$ applied to $|G_1\rangle\otimes|G_2\rangle$ leaves, up to a Pauli operator that
depends on the outcomes, the graph state of the graph on $(V_1\cup V_2)\setminus\{a,b\}$ whose edges are
the edges of $G_1$ and $G_2$ not incident to $a$ or $b$, together with every pair $\{u,w\}$ with
$u\in\nbr_{G_1}(a)$ and $w\in\nbr_{G_2}(b)$.
\end{lemma}

This is the standard graph rule for type-II fusion~\cite{Browne2005resource}. It is stated in this form in
Refs.~\cite{Zaidi2015near,Lee2023graph} and in full generality in Ref.~\cite{Lobl2025transforming}, and we
sketch the stabilizer calculation for completeness. The observable $X_aZ_b$ anticommutes exactly with $K_b$
and with $K_u$ for $u\in\nbr(a)$, and $Z_aX_b$ anticommutes exactly with $K_a$ and with $K_w$ for
$w\in\nbr(b)$. Measuring $X_aZ_b$ and then $Z_aX_b$ replaces $K_u$ by $K_uK_b$ and $K_w$ by $K_wK_a$.
Multiplying these by $Z_aX_b$ and $X_aZ_b$, respectively, removes their support on $a$ and $b$ and leaves
$X_u Z_{\nbr(u)\setminus\{a\}}Z_{\nbr(b)}$ and $X_w Z_{\nbr(w)\setminus\{b\}}Z_{\nbr(a)}$ up to sign, which
are the stabilizers of the graph in the lemma. Both of our stabilizer simulators also reproduce the rule on
random inputs (Sec.~\ref{sec:table}).

Fusing a vertex $a$ of a graph with a qubit of a fresh resource therefore adds one vertex, in one of three
ways that depend on the graph used to represent the resource. All three are available because
single-qubit Cliffords on a fresh resource are free. Fusing with a leaf of $P_3$ replaces $a$ by the center
of $P_3$ and attaches the other leaf to it, which up to relabeling adds a pendant vertex at $a$. Fusing
with the center of $P_3$ replaces $a$ by the two leaves, each with the neighborhood of $a$ and not adjacent
to each other, which adds a false twin. Fusing with a vertex of $K_3$ replaces $a$ by the other two
vertices, which are adjacent, and adds a true twin. We call these fusions merges.

\begin{proposition}[Counting bound]\label{prop:count}
Let a schedule with $g$ resources and $f$ fusions produce a connected graph on $n\ge 3$ vertices, and let
$H$ be the multigraph whose vertices are the resources, with one edge for each fusion joining the
resources that hold its two qubits. Then exactly one connected component $H_0$ of $H$ contains surviving
qubits. If $H_0$ has $g_0$ resources, $f_0$ fusions and cycle rank $t=f_0-g_0+1$, then $f\ge f_0=n-3+3t$.
In particular $\F(G)\ge n-3$, and every schedule has $f\equiv n\pmod 3$.
\end{proposition}

\begin{proof}
Fusions in different components of $H$ act on different tensor factors. If two components held surviving
qubits, the output would be a product state across them, which is impossible because every cut-rank of a
connected graph is positive, so that $|G\rangle$ is entangled across every bipartition. Counting the qubits
of $H_0$ gives $3g_0-2f_0=n$, and with $g_0=f_0-t+1$ this becomes $f_0=n-3+3t$, where $t\ge 0$ because $H_0$
is connected. The same count over all of $H$, whose other components have all their qubits fused, gives
$3g-2f=n$ and hence $f\equiv n\pmod 3$.
\end{proof}

The Lean theorem \code{fusion\_cost\_lower\_bound} proves a counting form of this bound in a merge model of
schedules that also tracks fusions which split or consume a component.

\section{The rank-width bound and the characterization}\label{sec:main}

A rank decomposition of a graph $G$ on at least two vertices is a tree $T$ whose internal nodes have degree
three and whose leaves are the vertices of $G$. Each edge $e$ of $T$ splits the leaves into two sets
$(A_e,V\setminus A_e)$, and the width of $T$ is $\max_e\rho_G(A_e)$. The rank-width $\rw(G)$ is the least
width of a rank decomposition of $G$~\cite{Oum2005rank,Oum2017rank}. A connected graph on at least two
vertices has $\rw(G)\ge 1$, and $\rw(G)\le 1$ if and only if $G$ is distance-hereditary~\cite{Oum2005rank}.

\begin{proof}[Proof of Theorem~\ref{thm:rw}]
Take a schedule with $f=\F(G)$ fusions that produces $G$, and let $H_0$, $g_0$, $f_0$ and $t$ be as in
Proposition~\ref{prop:count}. A single-qubit gate commutes with every operation that does not act on its
qubit, so each gate on a fused qubit may be applied to its resource before any fusion, and each gate on a
surviving qubit may be applied at the end, where it does not change Schmidt ranks. Fusions act on disjoint
pairs of qubits and commute with one another. Conditioned on the outcomes, the state of the surviving qubits
is then, up to normalization and single-qubit unitaries,
\[
  \psi=\Big(\bigotimes_{e}\langle\beta_e|\Big)\Big(\bigotimes_{r}|R_r\rangle\Big),
\]
where $r$ runs over the resources of $H_0$, $|R_r\rangle$ is resource $r$ with its gates applied, $e$ runs
over the fusions of $H_0$, and $\langle\beta_e|$ is the two-qubit bra of the outcome of $e$. The components of
$H$ without surviving qubits contribute only a scalar factor. Choose a spanning tree $S$ of $H_0$, so that
$t$ fusions of $H_0$ lie outside $S$, and form a tree $T$ from $S$ by attaching each surviving qubit as a leaf
to its resource. Each of the three qubits of a resource either survives, is fused along an edge of $S$, or
is fused along an edge outside $S$, so every resource has degree at most three in $T$. Repeatedly deleting
resource nodes of degree one and suppressing nodes of degree two turns $T$ into a rank decomposition $T'$
of $G$, once the surviving qubits are identified with the vertices of $G$.

Each edge of $T'$ arises from a path in $T$, and all edges of that path induce the same cut of the
surviving qubits. If the path is a single leaf edge, the cut separates one vertex and has cut-rank at most
one. Otherwise the path contains an edge $e$ of $S$. Removing $e$ splits the resources of $H_0$ into two
sets $X$ and $Y$, and the surviving qubits into $A$ and $B$, and the fusions joining $X$ and $Y$ are $e$
together with some of the $t$ fusions outside $S$. After contracting the fusions inside $X$ and inside $Y$,
each remaining bra $\langle\beta|$ is a two-qubit state and hence a sum of two product bras. Expanding them
writes $\psi$ as a sum of at most $2^{1+t}$ products of a state on $A$ and a state on $B$, so the Schmidt
rank of $\psi$ across $(A,B)$ is at most $2^{1+t}$. Since $\psi$ is $\LC$-equivalent to $|G\rangle$, it
follows that $\rho_G(A)\le 1+t$. Every edge of $T'$ therefore has cut-rank at most $1+t$, so
$\rw(G)\le 1+t$, and Proposition~\ref{prop:count} gives $\F(G)=f\ge n-3+3t\ge n-3+3(\rw(G)-1)$.
\end{proof}

The argument is the fusion analogue of bounding the Schmidt-rank width of a tensor network by its bond
structure~\cite{VandenNest2006universal,VandenNest2007classical}, and it is close in spirit to the bound
$\F\ge\mathrm{lrw}(G)-1$ in terms of the linear rank-width in Ref.~\cite{Lobl2025generating} and to the gate
bound of Ref.~\cite{Kumabe2026complexity}. Since only the cycle rank of the fusion network enters, the bound
holds for any destructive two-qubit measurement with pure outcomes, of which Bell-type fusions are a special
case.

\begin{proof}[Proof of Theorem~\ref{thm:main}]
($\Rightarrow$) If $\F(G)=n-3$, then $\rw(G)\le 1$ by Theorem~\ref{thm:rw}, so $G$ is
distance-hereditary~\cite{Oum2005rank}.

($\Leftarrow$) A connected distance-hereditary graph on $n\ge 3$ vertices is obtained from a connected
graph on three vertices, $P_3$ or $K_3$, by $n-3$ additions of pendant vertices, false twins and true twins,
with every intermediate graph connected~\cite{Bandelt1986distance}. Start from one resource represented as
$P_3$ or $K_3$, and realize each addition at a vertex $a$ by a merge with a fresh resource, at a leaf of
$P_3$, at the center of $P_3$ or at a vertex of $K_3$. By Lemma~\ref{lem:merge} this schedule produces $G$
with $n-3$ fusions, and Proposition~\ref{prop:count} gives $\F(G)=n-3$.
\end{proof}

At the minimum, the proof of Theorem~\ref{thm:rw} shows that $t=0$, so the fusion network is a tree on
$n-2$ resources whose $n-3$ fusions each join two components. The constructive direction can also be
obtained from the split decomposition. Every node of the split decomposition of a distance-hereditary graph
is a star or a complete graph~\cite{Bouchet1988transforming,Hammer1990completely,Gioan2012split}. Connolly
\emph{et al.}~\cite{Connolly2026efficient} assemble the graph state from the graph states of the nodes by
type-II fusions along the split tree, and Lee and Jeong~\cite{Lee2023graph} build an $m$-vertex star from
$m-2$ GHZ states with $m-3$ fusions. Together these again use $n-2$ resources and $n-3$ fusions.

\section{Formal verification}\label{sec:lean}

The Lean development is in the directory \code{lean/EmpiricistLean} of the accompanying repository and
uses \textsc{mathlib} at release \code{v4.31.0}. The ancillary files contain the modules behind every
declaration cited here as a self-contained Lake project, together with a script that builds it and prints
the axioms of each declaration (Table~\ref{tab:lean} in Appendix~\ref{app:lean}). All of them depend only
on the axioms \code{propext}, \code{Classical.choice} and \code{Quot.sound}. The development covers the
combinatorics of the constructive direction of Theorem~\ref{thm:main}. The steps that involve quantum
states, Lemma~\ref{lem:merge} and the proof of Theorem~\ref{thm:rw}, are not formalized.

Graphs are \textsc{mathlib} \code{SimpleGraph}s on an arbitrary vertex type, and the production predicates
are defined on concrete graphs. Let $\Prod^{+}_k(G)$ state that $G$ is obtained from
$P_3$ by $k$ steps, each of which adds a pendant vertex (the leaf-merge, \code{ghz3LeafMerge}) or a false
twin (the center-merge), interleaved with any number of local complementations and isomorphisms (the
inductive \code{ProducibleByExt}). The lemma \code{ghz3CenterMerge\_iso\_addFalseTwin} identifies the
center-merge, written as the graph rewrite of Lemma~\ref{lem:merge}, with false-twin addition up to
isomorphism. Let $\DHclass$ be the class generated from $P_3$ and $K_3$ by pendant, false-twin and
true-twin additions, local complementation and isomorphism (the inductive \code{PendantTwinBuildable}).
Without the local-complementation constructor these are the connected distance-hereditary graphs on at
least three vertices~\cite{Bandelt1986distance}, and the constructor adds no graphs, since that class is
closed under local complementation~\cite{Oum2005rank}. This identification is not formalized.

\begin{theorem}[Model characterization, Lean]\label{thm:model}
For a finite graph $G$, $\Prod^{+}_{|V(G)|-3}(G)$ holds if and only if $G\in\DHclass$
(\code{dh\_characterization}).
\end{theorem}

The forward implication is a structural induction on $\Prod^{+}$, each of whose constructors is also a
constructor of $\DHclass$. The converse (\code{dh\_forward}) is a structural induction on $\DHclass$. A
pendant addition is a leaf-merge and a false-twin addition is a center-merge, and a true-twin addition is
realized by a local complementation, a leaf-merge and a second local complementation, through the identity
\begin{equation}\label{eq:truetwin}
  \tau_u\big(\mathrm{addPendant}(G,u)\big)\;=\;\mathrm{addTrueTwin}\big(\tau_u(G),\,u\big)
\end{equation}
(\code{addPendant\_localComplement\_center}) and the fact that $\tau_u$ is an involution. Combined with
\code{fusion\_cost\_lower\_bound}, the theorem \code{dh\_min\_fusions} states the constructive direction of
Theorem~\ref{thm:main} at the level of the model. A further theorem, \code{floor\_schedule\_iff\_dh}, extends
Theorem~\ref{thm:model} to a schedule model with one additional constructor, which may replace the current
graph by an arbitrary graph with two fewer vertices at the cost of one fusion. This over-approximates a
fusion of two qubits of the same component, and at the minimum a count of vertices excludes it.

Before the general theorem, the development treats several families separately. For each family, the
Lean theorem states both the general lower bound and the existence of a construction that attains it.
Paths, stars and all trees are handled by induction on the removal of leaves, $K_n$ is one local complementation away
from a star, and $K_{m,n}$ is three local complementations away from a double star. All of these graphs are
distance-hereditary.

\begin{remark}\label{rem:concrete}
Although every statement about $\F$ is invariant under local complementation, the inductions cannot be
carried out on $\sim_{\LC}$ classes, because pendant addition does not descend to them. The graphs $P_3$
and $K_3$ are $\LC$-equivalent. A pendant at the center of $P_3$ gives the four-vertex star, whereas a
pendant at a vertex of $K_3$ gives the paw, which is $\LC$-equivalent to the path $P_4$ and lies in a
different orbit from the star. The development therefore carries concrete graphs through every induction
and applies local complementation as an explicit step. An earlier production model contained only the leaf-merge and
therefore generated only trees. It reaches the minimum only on the local-complementation classes of trees,
which make up $93$ of the $175$ distance-hereditary orbits with $3\le n\le 9$, one for each tree on $3$ to
$9$ vertices, since locally equivalent trees are isomorphic~\cite{Bouchet1988transforming}. Adding the
center-merge removed this restriction.
\end{remark}

\section{Minimum fusion counts for small graphs}\label{sec:table}

We assigned every connected graph with $3\le n\le 9$ to its local-complementation orbit. There are
$1$, $2$, $4$, $11$, $26$, $101$ and $440$ orbits for $n=3,\dots,9$, or $585$ in total, in agreement with Danielsen and
Parker~\cite{Danielsen2006classification}. By Proposition~\ref{prop:count} the search proceeds in tiers
$\F=n-3$, $\F=n$, and so on. Tier~0 is a breadth-first search over graphs in which local complementation is
free and a merge at any vertex, with a fresh resource fused at a leaf or at the center, costs one fusion;
true twins arise through Eq.~\eqref{eq:truetwin}. Tier~1 grows graphs by merges to $n+2$ vertices and then
performs one fusion of two qubits of the same component. Both tiers are exhaustive by the normal form of
Appendix~\ref{app:normal}, which shows that any schedule can be reordered so that the merges come first
and grow a single component, and that single-qubit Cliffords on the qubits of a merge can be absorbed by
local complementations. Every fusion in the search is computed in the fixed basis of Sec.~\ref{sec:model}.
For the single fusion inside the component in Tier~1 we repeated the search with each of the six classes of
single-qubit Clifford on one of the two fused qubits, which together with the symmetry of the measured pair
covers all choices on both qubits, and for $n\le 5$ we tried all $36$ pairs. The additional Clifford
choices did not change the set of orbits reached by Tier~1.

Each recorded construction is replayed by two stabilizer simulators with independent implementations of
fusion and of graph extraction, a tableau simulator built on \textsc{stim}~\cite{Gidney2021stim} and an
implementation over $\mathrm{GF}(2)$. The two share the code that replays a construction, the graph-level
local complementation and the canonicalization of orbits, which uses \textsc{nauty}~\cite{McKay2014practical}.
The closed-form merge rule used in the inner loop of the search agrees with both simulators on randomized
inputs. The resulting table, a tablebase in the terminology of
Ref.~\cite{Lobl2025generating}, has one row for each of the $585$ orbits, with an orbit representative and
either an exact value with a witness schedule or a lower bound, and it is stored as a single file
identified by a hash of its contents. Nine exact values above the minimum were added later from constructions proposed
by the language model.

\begin{table}[t]
\centering
\begin{tabular}{lrrrrrrrr}
\toprule
$n$ & 3 & 4 & 5 & 6 & 7 & 8 & 9 & total\\
\midrule
connected orbits               & 1 & 2 & 4 & 11 & 26 & 101 & 440 & 585\\
$\F=n-3$ (distance-hereditary) & 1 & 2 & 3 &  8 & 15 &  42 & 104 & 175\\
$\F=n$                         & 0 & 0 & 1 &  2 &  7 &   6 &   0 &  16\\
$\F=n+3$                       & 0 & 0 & 0 &  1 &  2 &   0 &   0 &   3\\
open, $\F\ge n$                & 0 & 0 & 0 &  0 &  0 &  52 & 312 & 364\\
open, $\F\ge n+3$              & 0 & 0 & 0 &  0 &  2 &   1 &  24 &  27\\
\bottomrule
\end{tabular}
\caption{Connected local-complementation orbits by number of vertices. The orbits with $\F=n-3$ are
exactly the distance-hereditary ones. All other orbits have rank-width two, except for the $25$ open orbits
with $n\in\{8,9\}$ and rank-width three, whose bound $\F\ge n+3$ comes from Theorem~\ref{thm:rw}. The two
open orbits with $n=7$ have $\F\ge n+3$ from the exhaustive Tier~1 search.}
\label{tab:counts}
\end{table}

Table~\ref{tab:counts} gives the number of orbits in each class. An independent pendant-and-twin
elimination~\cite{Bandelt1986distance} on all $585$ representatives confirms that the orbits with $\F=n-3$ are exactly the distance-hereditary
ones. Nineteen orbits have exact values above the minimum. For the sixteen orbits with $\F=n$, the upper
bounds come from constructions and the lower bounds from Theorem~\ref{thm:main}. The three orbits with
$\F=n+3$ are a six-vertex orbit with $\F=9$, realized with eight resources, nine fusions and three local
complementations, and two seven-vertex orbits with $\F=10$. Their lower bounds come from the exhaustive
Tier~1 search.

To test the rank-width bound, we computed the rank-width of every orbit above the minimum exactly, by
enumerating all rank decompositions. Every orbit with $n\le 7$ that is not distance-hereditary has
rank-width two, and so do all sixteen orbits with $\F=n$, which therefore attain the bound $\F\ge n$ of
Theorem~\ref{thm:rw}. The three orbits with $\F=n+3$ also have rank-width two, and the bound falls short of
their cost by three. Above rank-width one, therefore, the fusion cost is not a function of rank-width. For
$25$ open orbits of rank-width three, one with $n=8$ and $24$ with $n=9$, Theorem~\ref{thm:rw} strengthens
the bound $\F\ge n$ established by the search to $\F\ge n+3$.

\section{Computational methods}\label{sec:method}

The search and the formal proofs were produced with Empiricist, an open-source system in which a language
model proposes conjectures, constructions and complete Lean modules as structured output. The model never
executes code, and the system runs every check in a sandbox. It records each claim as heuristic,
conjectured, verified up to a given size, certified or formalized, in increasing order of strength, or as
refuted, and a claim advances only when a verifier accepts evidence for it. A claim supported only by the
model's assertion is recorded at most as heuristic, and the status refuted is permanent. Each verifier
must pass a fixed suite of test cases before use, and artifacts and verifiers are identified by content
hashes, a verifier by the hashes of its code and of its test suite. Fusion outcomes are computed by the two simulators described in Sec.~\ref{sec:table}, and a
construction is accepted only when both agree on the resulting state up to local Clifford equivalence.
Lean modules are elaborated in a sandbox without network or process creation, against a frozen,
read-only copy of \textsc{mathlib}. The kernel result is checked again, the axioms must lie within
$\{\code{propext},\code{Classical.choice},\code{Quot.sound}\}$, every import must belong to a fixed trusted
set, and the build directory is cleared after elaboration. Before it was used, the checker was revised in
response to adversarial submissions, including tampered compiled files, substitution of files between
checking and use, and elaborators that hide axioms.

The distance-hereditary pattern was first noticed by the coding assistant used to build the system, when
it examined the table of all $585$ orbits. Working within the system, the model was then given only the
$44$ orbits with $3\le n\le 7$ as edge lists, each marked according to whether it attains $\F=n-3$, with no invariant or class name and with only the information
that the property is invariant under local complementation. It proposed that the orbits attaining the
minimum are the distance-hereditary graphs, recalled the standard equivalent characterizations, and
identified the orbit of $C_5$ as the unique orbit on five vertices that does not attain it. The proposal was
recorded as a conjecture, and a deterministic test for distance-heredity on all $585$ orbits, including the
$541$ orbits on eight and nine vertices that the model had not seen, agreed with it in every case. The
model then wrote the Lean proofs in a loop in which the checker returns the remaining proof obligations and
the model revises the module, without running Lean itself. Two errors were caught during the work.
The first statement of a family theorem had its construction built into the claim. Such a statement is
well typed, so the Lean checker cannot detect it, and it was caught when the statement was reviewed. While
preparing this paper, we found by an exact computation of rank-widths that a note had given the linear
rank-width of the six-vertex orbit with $\F=9$ as its rank-width. Related uses of machine learning in mathematics include conjectures guided by
data~\cite{Davies2021advancing}, new constructions in combinatorics and analysis found by program search with
language models~\cite{RomeraParedes2024mathematical,Novikov2025alphaevolve,Georgiev2025mathematical}, and
machine-generated proofs checked by a symbolic deduction engine~\cite{Trinh2024solving} or by the Lean
kernel~\cite{Hubert2026olympiad}.

\section{Discussion}\label{sec:discussion}

Extending the formal verification from the combinatorial core of the constructive direction of
Theorem~\ref{thm:main} to Lemma~\ref{lem:merge} and to the proof of Theorem~\ref{thm:rw} would require
definitions of graph states, Pauli measurements and schedules in Lean, and rank-width is not yet
available in \textsc{mathlib}.

The data of Sec.~\ref{sec:table} show that $\F$ is not a function of rank-width beyond rank-width one, and
it is not known which invariant separates $\F=n$ from $\F=n+3$ among graphs of rank-width two.
Theorem~\ref{thm:main} gives a polynomial-time test for $\F(G)=n-3$, since distance-hereditary graphs are
recognized by pendant-and-twin elimination~\cite{Bandelt1986distance}, but the complexity of computing $\F$
in general is open, as is the question whether analogous characterizations hold for larger GHZ states or
for linear cluster states as resources.

\section*{Data and code availability}

The repository \url{https://github.com/sghowell/empiricist} (MIT license) contains the system, the Lean
development, the table of minimum fusion counts and the record of claims. The ancillary files of this
paper contain a self-contained Lake project with the Lean modules for all declarations in
Table~\ref{tab:lean} and a script that builds the modules and prints their axioms. They also contain the
table, the nine later constructions and the scripts used to check them. The scripts recompute every number
in Sec.~\ref{sec:table}, including the exact rank-widths, replay the constructions with both simulators and
repeat Tier~1 with single-qubit Cliffords on the fused qubits. A README maps each statement of the paper to a file, and a SHA-256 manifest lists every
file.

\begin{acknowledgments}
This work was carried out with AI systems working under my direction. Claude Code (Anthropic, model
Claude Opus~4.8) built the first version of Empiricist, including its exhaustive search, and first noticed
the distance-hereditary pattern in the resulting table. Within Empiricist, the language model Fable~5
(Anthropic) proposed the nine later constructions, re-derived the distance-hereditary characterization from
the orbits with $3\le n\le 7$ and wrote the Lean proofs, emitting only machine-checkable output and never
executing code. Claude Code (Anthropic, model Claude Opus~5.5) surveyed the literature, proved
Theorem~\ref{thm:rw}, wrote the remaining proofs, recomputed every number from the committed data and
drafted the manuscript. Codex (OpenAI, model gpt-6-astra) independently checked the proofs, the claims and
the relation to previous work; the errors it found have been corrected. I chose the problem and take full
responsibility for the content of this paper.
\end{acknowledgments}

\appendix

\section{The normal form behind the search}\label{app:normal}

The exhaustiveness of Tiers~0 and~1 rests on two lemmas.

\begin{lemma}[Absorbing single-qubit Cliffords]\label{lem:frame}
Let $v$ be a vertex of a graph $G$ with at least one neighbor, and let $C$ be a single-qubit Clifford
unitary. Then there are a graph $G'$, obtained from $G$ by local complementations, and a product $D$ of
single-qubit Cliffords acting trivially on $v$, such that $C_v|G\rangle=D\,|G'\rangle$ up to a Pauli
operator and a global phase.
\end{lemma}

\begin{proof}
Single-qubit Clifford unitaries modulo Pauli operators and phases form a group isomorphic to the symmetric
group on $\{X,Y,Z\}$, acting by conjugation, and $e^{-i\pi X/4}$ and $e^{i\pi Z/4}$ act as the transpositions
$(Y\,Z)$ and $(X\,Y)$, which generate it. By Eq.~\eqref{eq:lc}, local complementation at $v$ multiplies the
local unitary on $v$ by $e^{-i\pi X_v/4}$, and local complementation at a neighbor $u$ of $v$ multiplies it
by $e^{i\pi Z_v/4}$. Neither operation removes the edge $\{u,v\}$, since $\tau_v$ changes only edges inside
$\nbr(v)$ and $\tau_u$ only edges inside $\nbr(u)$. A suitable word in $\tau_v$ and $\tau_u$ therefore gives
a graph $G'$ and a product $U=\bigotimes_w U_w$ of single-qubit Cliffords with $U|G\rangle=|G'\rangle$ up to a
Pauli operator and $U_v=C$ modulo Paulis. Then $C_v|G\rangle=\big(\bigotimes_{w\neq v}U_w^{-1}\big)U|G\rangle$,
which equals $D|G'\rangle$ with $D=\bigotimes_{w\ne v}U_w^{-1}$ up to a Pauli operator.
\end{proof}

\begin{lemma}[Single-component normal form]\label{lem:normal}
Suppose a schedule with $f$ fusions produces a connected graph $G$ on $n$ vertices. Then $G$ is also
produced by a schedule with $f'\le f$ fusions and $g'$ resources, $3g'-2f'=n$, of the following form. One
resource serves as the seed, and each of the other $g'-1$ resources is attached by a merge, a fusion between
a qubit of the current component, which is a connected graph state, and a qubit of the fresh resource, with
no single-qubit Clifford on either fused qubit. Local complementations of the component may occur between
merges. The remaining $f'-g'+1$ fusions follow all merges and act on two qubits of the component, and
single-qubit Cliffords on the surviving qubits come last.
\end{lemma}

\begin{proof}
A single-qubit Clifford on a qubit $q$ commutes with every operation that does not act on $q$, and the only
other operation that acts on $q$ is the fusion that consumes it, if there is one. We first move each such
gate to just after the preparation of its resource, to just before the fusion that consumes $q$, or to the
end of the schedule if $q$ survives.

Next we delete the components of the multigraph $H$ of Proposition~\ref{prop:count} that contain no
surviving qubit, since they act on separate tensor factors. The resulting $H$ is connected, with $g'$
resources and $f'\le f$ fusions, and $3g'-2f'=n$.

We then reorder the fusions. Choose a spanning tree of $H$, a root resource, and a breadth-first order of
the resources from the root along the tree. For each resource $c$ other than the root, with parent $p$,
pick one fusion joining $p$ and $c$, and call these fusions the merges. We place the merges first, in the
order of their child resources, followed by all other fusions in their original relative order. Fusions act
on disjoint pairs of qubits and commute, so the final state on the surviving qubits, conditioned on the
outcomes, does not change, and only the Pauli corrections may differ. When the merge of $c$ is performed,
its qubit in $p$ has not been consumed and $c$ has not yet been touched, so the merge fuses a qubit of the
component grown so far with a qubit of a fresh resource. Every later fusion acts on two qubits of that
component.

Finally we remove the Cliffords on the qubits of the merges, processing the merges in order. Before the merge
of $c$, the component is the graph state of a connected graph $B$ on at least three vertices, up to
single-qubit Cliffords, because the seed has three vertices and each merge adds one vertex and keeps the graph
connected by Lemma~\ref{lem:merge}, both joined neighborhoods being nonempty. The fused qubit $a$ of the
component has a neighbor, so Lemma~\ref{lem:frame} replaces the Clifford on $a$ by local complementations of
$B$ and Cliffords on other qubits, which are carried forward. On the side of the fresh resource,
Lemma~\ref{lem:frame} applied to its three-vertex graph removes the Clifford on the fused qubit and leaves
$P_3$ with some center, or $K_3$. The merge is then a fusion without single-qubit Cliffords on either fused
qubit. Cliffords carried onto qubits that later merges consume are removed in the same way when those merges
are reached, Cliffords on qubits consumed by the final fusions inside the component remain, as the lemma
allows, and the rest move to the end.
\end{proof}

At the minimum the normal form has no fusion inside the component, so Tier~0 is exhaustive for $\F=n-3$.
Tier~1, which applies one such fusion after all merges, is exhaustive for $\F=n$, given the repetition of the
search with single-qubit Cliffords on its two fused qubits described in Sec.~\ref{sec:table}.

\section{Principal Lean declarations}\label{app:lean}

\begin{table}[h]
\centering
\begin{tabular}{lll}
\toprule
declaration & statement & file\\
\midrule
\code{fusion\_cost\_lower\_bound} & $f\ge n-3$ in a merge model of schedules & \code{Foundation.lean}\\
\code{tree\_min\_fusions} & construction at the minimum for every tree & \code{TreeThm.lean}\\
\code{pathGraph\_min\_fusions} & paths & \code{Paper/PathGraph.lean}\\
\code{star\_min\_fusions} & stars & \code{Paper/Star.lean}\\
\code{complete\_min\_fusions} & complete graphs $K_n$ & \code{Paper/Complete.lean}\\
\code{completeBipartite\_min\_fusions} & complete bipartite graphs $K_{m,n}$ & \code{Paper/CompleteBipartite.lean}\\
\code{ghz3CenterMerge\_iso\_addFalseTwin} & the center-merge is false-twin addition & \code{CenterMerge.lean}\\
\code{addPendant\_localComplement\_center} & the identity~\eqref{eq:truetwin} & \code{TrueTwin.lean}\\
\code{dh\_forward}, \code{dh\_min\_fusions} & $\DHclass\Rightarrow\Prod^{+}_{n-3}$ & \code{DHCharacterization.lean}\\
\code{dh\_characterization} & $\Prod^{+}_{n-3}(G)\Leftrightarrow G\in\DHclass$ & \code{DHCharacterization.lean}\\
\code{floor\_schedule\_iff\_dh} & the over-approximating schedule model & \code{Paper/FloorBridge.lean}\\
\bottomrule
\end{tabular}
\caption{Principal Lean declarations. Files without a directory are library modules of the repository, and
files in \code{Paper/} are modules recorded as checked artifacts in its record of claims.}
\label{tab:lean}
\end{table}

\bibliography{refs}

\end{document}
